% !TEX root = userguide.tex

\def\headdate{6th December 2023}

\section{Viscoelastic fluid flow using the b-tensor formulation}
\label{sec:CDTformulation}

(These examples in this section have been developed by Mick Carrozza and Michelle Spanjaards of the TU/e)

\subsection{Theory}

In the b-tensor formulation the conformation tensor $\ten c$ is decomposed
according to $\ten c=\ten b\cdot\ten b^T$. The differential form for $\ten b$ is similar to the deformation gradient, i.e.\ with just one $\ten L\cdot\ten b$ term instead of two in the conformation tensor formulation. In \tfem\ four different variants are available:
\begin{enumerate}
\item  The contravariant deformation tensor (CDT)
formulation introduced in \cite{hutter_fluctuating_2018}. Details on the
numerical implementation of the CDT-formulation can be found in \cite{carrozza_viscoelastic_2019}. 
In \tfem\ this is called \texttt{bvariant=1}.
\item The symmetric variant of Balci et al.\ \cite{balci_symmetric_2011}.
In \tfem\ this is called \texttt{bvariant=2}.
\item The Cholesky variant of Vaithianathan and Collins \cite{vaithianathan_numerical_2003}.
In \tfem\ this is called \texttt{bvariant=3}.
\item The Cholesky variant with
log transformation of the diagonal components of Vaithianathan and Collins \cite{vaithianathan_numerical_2003} .
In \tfem\ this is called \texttt{bvariant=4}.
\end{enumerate}
The first variant is the most simple one (no additional velocity gradient terms), but requires rotation reinitialization and the solution for the $\ten b$-tensor is generally unsteady in steady flows.
The other variants do not have these two drawbacks (rotation reinitialization and unsteadiness), but an additional term $\ten b\cdot\ten A$ is required, where the anti-symmetric tensor
$\ten A$ depends on $\ten b$ and $\ten L$. Early testing shows that the CDT formulation (\texttt{bvariant=1}) is the most stable one in our current implementation, similar to the log-conformation representation. Therefore it is recommended to use \texttt{bvariant=2,3,4} only for comparison and not for production runs.

In the examples that follow, except the ones with thermal fluctuations, all variants can be tested.

In the CDT-formulation of Carrozza et al.\ \cite{carrozza_viscoelastic_2019} it is required to do so-called rotation reinitialization, which basically sets the rotation tensor $\ten R$ of $\ten b_n$ back to the unity tensor. For multi-step schemes, values of $\ten b$ at older times obtain the same additional rotation as $\ten b_n$. In Carrozza et al.\ \cite{carrozza_viscoelastic_2019} the rotation reinitialization is performed \emph{after} the time step for all the nodal values with $\ten b_{n+1}$ as a reference (and also values at older time steps, if needed). This means that $\ten b_{n+1}=\sqrt{\ten c_{n+1}}$ is a symmetric tensor after the time step. Note, that $\ten b_{n+1}$ becomes $\ten b_n$ in the next time step. Alternatively, rotation reinitialization can be performed on element level in each integration point \emph{at the start} of the time step with $\ten b_n$ as a reference. This means that $\ten b_{n+1}$ will be an unsymmetric tensor after solving the CDT system! Also note, that the reinitialization is performed on the element values only and not on the system vector(s) $\ten b_n$, ($\ten b_{n-1}$). Hence the $\ten b$ tensors between the two types of reinitialization cannot be directly compared. Only the $\ten c$ tensors should be comparable.
Element level rotation reinitialization simplifies the main program, removing the burden of reinitialization at the main level. However, the implementation has been limited to BDF-type time integration schemes (no trapoidal rule, for example).

\subsection{Planar Couette flow of an upper-convected Maxwell fluid (UCM) for a
 Weissenberg number (Wi) of 10}
\label{sec:couette_b}

We consider an upper-convected Maxwell fluid model (UCM) for a planar Couette problem on a unit square
(see Fig.~\ref{fig:couette}). On the upper boundary a tangential
velocity of 1 and a normal velocity of zero is imposed.
On the lower wall the velocity components are
zero. In horizontal direction periodical boundary conditions are assumed.
In the lower left corner the pressure value is set to zero (Dirichlet).

At the start of the flow the b-tensor field will be set to the exact steady state value
of the square root of the conformation tensor for variants 1 and 2. For the variants 3 and 4 the initial value of the b-tensor is set to the lower diagonal matrix of the Cholesky decomposition 
of the exact steady state value of the conformation tensor. In addition, for variant 4 the log is taken of the diagonal components of the Cholesky lower triangular matrix.

Since the contravariant deformation tensor field (in variant 1)
 is not a steady field there is a small time discretization error 
even in a steady state shear flow. Therefore, first the reference steady state solution is found by performing a number of time steps. Then, a small random perturbation is added to the contravariant deformation tensor. 

The example file is \texttt{couette15.f90} (second-order stress-implicit
scheme) and can be obtained by typing at the
command line:
\medskip
\begin{quote}
   \texttt{getexample couette}
\end{quote}
%
The example can be run by
\begin{quote}
   \texttt{couette15 < input\_couette15.txt}
\end{quote}
%
Time-dependent stability is similar to the standard formulation using the conformation tensor.

Example \texttt{couette15a.f90} is similar to example \texttt{couette15.f90}, but uses element level rotation reinitialization instead of the nodal based one.

\subsection{Planar Couette flow of an Oldroyd-B fluid for a
 Weissenberg number (Wi) of 10 (second-order stress-explicit scheme)}
\label{sec:couette_b_OldB}

The example file is \texttt{couette16.f90} and can be obtained by typing at the
command line:
\begin{quote}
   \texttt{getexample couette}
\end{quote}
%
The example can be run by
\begin{quote}
   \texttt{couette16 < input\_couette16.txt}
\end{quote}

Example \texttt{couette16a.f90} is similar to example \texttt{couette16.f90}, but uses element level rotation reinitialization instead of the nodal based one.

\subsection{Flow around a cylinder between two plates}
\label{sec:confcyl1}
The examples can be obtained by typing at the command line:
\begin{quote}
   \texttt{getexample confined\_cylinder\_b}
\end{quote}

Examples \texttt{cylinder1.f90} and \texttt{cylinder2.f90} describe the flow
around a fixed cylinder positioned between two walls. The numerical scheme is
stress-explicit (see Sec.~\ref{sec:explicitstress}) and stress-implicit (see
Sec.~\ref{sec:implicitstress}) for the two examples, respectively.

Examples \texttt{cylinder1a.f90} and \texttt{cylinder2a.f90} are similar to examples \texttt{cylinder1.f90} and \texttt{cylinder2.f90}, respectively, but use element level rotation reinitialization instead of the nodal based one.

A default coarse mesh (\texttt{confined\_cylinder1.msh}) is supplied, but in the
subdirectory \texttt{meshes} files are available to generate more refined 
meshes.

For comparison the examples \texttt{cylinder\_c1.f90} and \texttt{cylinder\_c2.f90} are also available.
These are similar to \texttt{cylinder1.f90} and \texttt{cylinder2.f90}, respectively, however 
the problems are solved using the (log-)conformation tensor. By considering the difference between these files, users are able to transfer their problems to using the b-formulation instead.

\subsection{Extrudate swell}

The examples in this section are the extrudate swell examples in \cite{hulsen_numerical_2021}.
The examples can be obtained by typing at the command line:
\begin{quote}
   \texttt{getexample extrudate\_swell\_b}
\end{quote}

Examples \texttt{extrudate\_swell2d\_b.f90} and \texttt{extrudate\_swell3d\_b.f90} describe extrudate swell 
of a viscoelastic fluid flowing from a die for a 2D planar and a 3D cylindrical geometry, respectively.

Examples \texttt{extrudate\_swell2d\_ba.f90} and \texttt{extrudate\_swell3d\_ba.f90} are similar to examples \texttt{extrudate\_swell2d\_b.f90} and \texttt{extrudate\_swell3d\_b.f90}, respectively, but use element level rotation reinitialization instead of the nodal based one.

For comparison the examples \texttt{extrudate\_swell2d\_c.f90} and \texttt{extrudate\_swell3d\_c.f90} are also available.
These are similar to \texttt{extrudate\_swell2d\_b.f90} and \texttt{extrudate\_swell3d\_b.f90}, respectively, however 
the problems are solved using the (log-)conformation tensor. 

\subsection{Flow in curved channel with a square cross section (Dean flow, axisymmetric with swirl)} 
\label{sec:Dean2Dimplbten}

This problem is in file \texttt{channel14.f90} and is similar to the problem \texttt{channel11.f90} solved in Sec.~\ref{sec:Dean2Dimpl}. 

Example \texttt{channel14a.f90} is similar to example \texttt{channel14.f90}, but uses element level rotation reinitialization instead of the nodal based one.

\subsection{Equilibrium and homogeneous flows with fluctuations}
\label{sec:startup_b_fluct}

The stochastic differential equation for viscoelastic fluids with fluctuations, 
written in the CDT-formulation, is integrated in time while imposing a constant 
velocity gradient. Forward Euler time stepping is used. The initial condition 
for the contravariant deformation can for instance be chosen to be the 
equilibrium solution for the deterministic model, or the square root of the
average conformation in equilibrium for the stochastic UCM model. Single-point
computations in 3D are performed for equilibrium and shear or uniaxial
extensional flow. Pseudorandom numbers from a standard normal distribution are
drawn with a random-number generator using the Ziggurat method, 
see Sec.~\ref{sec:ziggurat}. A small enough time step should be used to prevent 
the solution to diverge due to finite-amplitude fluctuations.

The example file is \texttt{startup14.f90} and can be obtained by typing at the
command line:
\medskip
\begin{quote}
   \texttt{getexample viscoelastic\_models}
\end{quote}

